% Bản nháp chương 2 theo dàn bài mới, viết với bộ mẫu ở style/mau/.
% Chưa nối vào report.tex.
% ===========================================================================
\chapter{Chuẩn bị dữ liệu}
\label{sec:data}

Phần xử lý dữ liệu giữ ở mức tối thiểu, vì trọng tâm của bài tập là tối ưu hóa.
Nhóm chuyển các cột dạng chuỗi có đơn vị về số, mã hóa one-hot toàn bộ biến
định tính, bổ sung tích từng cặp và bình phương của các cột định lượng, rồi
chuẩn hóa mọi cột về trung bình 0 và độ lệch chuẩn 1. Các cột hằng bị loại vì
không mang thông tin, còn bình phương của 7 cột nhị phân bị loại vì nó tái tạo
đúng cột gốc.

Ba quyết định trong số đó chạm thẳng vào bài toán tối ưu hóa, và chương này đọc
cả ba theo cùng một trục: mỗi quyết định làm bài toán dễ đi hay khó đi bao
nhiêu, đo bằng số điều kiện $\kappa$ ở \eqref{eq:constants}. Ba mục sau lần lượt
xét cách tạo biến tương tác, cách chuẩn hóa cột, và cách chọn hệ số hiệu chỉnh.

\section{Ma trận thiết kế}
\label{sec:design-matrix}

Số thuộc tính sau khi mã hóa là $d = 280$. Trong 280 cột đó có 21 cột định
lượng gốc, 210 tích từng cặp, 14 bình phương và 35 cột one-hot, nên phần lớn
kích thước bài toán do biến tương tác sinh ra chứ không do biến định tính. Tích
từng cặp của 21 cột cho $\binom{21}{2} = 210$ cột, gấp sáu lần toàn bộ phần
one-hot của 6 biến định tính với 41 mức ở bảng~\ref{tab:dataset}. Tỉ lệ này sẽ
đảo nếu một biến định tính như \texttt{model} có hàng nghìn mức thay vì 23 mức,
khi đó khối one-hot mới là phần chi phối $d$.

\section{Tạo biến tương tác trên cột đã chuẩn hóa}
\label{sec:interaction}

Thứ tự giữa hai thao tác nhân cột và chuẩn hóa cột quyết định độ khó của bài
toán, dù không gian cột sinh ra thì như nhau. Lấy tích trực tiếp trên thang
gốc, tức nhân \texttt{original\_price} cỡ $10^{5}$ với
\texttt{screen\_size\_inches} cỡ 6, thì ma trận thiết kế thu được có
$L = \num{17.14}$ và 48 trị riêng dưới $10^{-6}$. Lấy tích trên các cột đã
chuẩn hóa thì $L = \num{2.68}$ và số trị riêng dưới $10^{-6}$ trở về đúng 7,
bằng mức của khối cột gốc.

Chênh lệch 6{,}4 lần về số điều kiện giữa hai cách hoàn toàn do cách biểu diễn.
Một cột tích trên thang gốc có phương sai lớn hơn các cột còn lại nhiều bậc nên
nó kéo riêng trị riêng lớn nhất lên, trong khi không gian cột thì không đổi.
Hai cách chỉ cho cùng một $\kappa$ khi mọi cột định lượng vốn đã cùng một thang
đo, và đó không phải trường hợp của bộ dữ liệu này.

\section{Chuẩn hóa cột và số điều kiện}
\label{sec:normalization}

\resultfig{normalization_iter}
  {Ảnh hưởng của việc chuẩn hóa cột lên tốc độ hội tụ của gradient descent, theo
   số vòng lặp. Trục tung là $f(w_k) - \fstar$ trên thang logarit.}
  {fig:normalization-iter}
\resultfig{normalization_time}
  {Ảnh hưởng của việc chuẩn hóa cột lên tốc độ hội tụ của gradient descent, theo
   thời gian chạy.}
  {fig:normalization-time}

Với cùng dữ liệu và cùng $\lambda$, số điều kiện là $\kappa = \num{268.3}$ sau
khi chuẩn hóa cột nhưng lên tới $\kappa = \num{46540}$ khi các cột giữ thang đo
không đồng nhất, tức chênh nhau 173 lần. Chênh lệch đó hiện thẳng ra ở tốc độ
hội tụ trên hình~\ref{fig:normalization-iter}: bản đã chuẩn hóa chạm giới hạn
số học $10^{-16}$ sau 295 vòng lặp, còn bản giữ thang đo gốc chạy hết
500 vòng mà sai số mới xuống cỡ $10^{-2}$, tức vẫn cao hơn mức bản chuẩn hóa
đạt được trong mươi vòng đầu tiên. Hệ số co rút mỗi vòng lặp của gradient
descent có dạng $1 - \bigO(1/\kappa)$, nên $\kappa$ lớn hơn 173 lần thì phần
sai số cắt được sau mỗi vòng cũng nhỏ đi theo đúng tỉ lệ ấy.

Chi phí mỗi vòng lặp không phụ thuộc cách biểu diễn cột, khoảng
\SI{28}{\milli\second} cho cả hai, nên hình~\ref{fig:normalization-time} chỉ là
hình~\ref{fig:normalization-iter} vẽ lại trên một thang khác: bản chuẩn hóa
xuống tới giới hạn số học trong \SI{8.4}{\second} còn bản kia sau
\SI{13.9}{\second} vẫn dừng ở $10^{-2}$. Do đó chuẩn hóa cột không phải một bước
tiền xử lý về thống kê mà là bước quyết định bài toán tối ưu hóa có giải được
trong thời gian hợp lý hay không. Chú ý rằng khoản lợi ấy chỉ tính cho các
phương pháp bậc một, vì
Newton ở chương~\ref{sec:mechanisms} vẫn hội tụ sau đúng một vòng lặp dù
$\kappa$ bằng bao nhiêu.

\section{Chọn hệ số hiệu chỉnh}
\label{sec:choose-lambda}

\resultfig{lambda_cv}
  {Sai số cross-validation 5 fold theo $\lambda$. Dải xám là một sai số chuẩn
   quanh giá trị nhỏ nhất, đường chấm là điểm cực tiểu và đường gạch là giá trị
   được chọn.}
  {fig:lambda-cv}

Đường cong cross-validation 5 fold ở hình~\ref{fig:lambda-cv} cho sai số trung
bình giống nhau tới bốn chữ số thập phân với mọi $\lambda$ từ $10^{-6}$ đến
$10^{-3}$, tức phẳng suốt ba bậc thập phân, vì mức hiệu chỉnh nhỏ hơn nhiễu của
dữ liệu thì không đổi được sai số dự báo. Trong khi đó, theo
\eqref{eq:mu-equals-lambda}, số điều kiện giữa hai đầu khoảng ấy chênh nhau
1000 lần. Vậy $\lambda$ vẫn phải chọn cẩn thận, nhưng tiêu chí chọn phải nằm
ngoài cross-validation. Đường cong chỉ có cực tiểu rõ khi $n$ nhỏ hơn nhiều
hoặc $d$ gần $n$, và bài toán này không rơi vào trường hợp đó.

Cực tiểu của một đường cong phẳng như vậy chỉ phản ánh nhiễu của phép chia
fold, nên nhóm dùng quy tắc một sai số chuẩn: chọn $\lambda$ lớn nhất mà sai số
cross-validation của nó vẫn nằm trong một sai số chuẩn của giá trị tốt nhất.
Quy tắc này cho $\lambda = \num{0.01}$ ứng với $\kappa = \num{268.3}$, trong
khi điểm cực tiểu nằm ở $\lambda = 10^{-4}$ ứng với $\kappa = \num{26831}$, tức
chênh nhau đúng 100 lần theo tỉ lệ giữa hai giá trị $\lambda$.

Cái giá phải trả về mặt thống kê rất nhỏ. RMSE trên tập kiểm tra là
\num{0.20601} tại $\lambda = \num{0.01}$ so với \num{0.20579} tại điểm cực tiểu
cross-validation, kém đi 0{,}1\%, và quy theo \eqref{eq:rmse-to-price} thì sai
số giá đi từ \num{4240} lên \num{4245} đơn vị tiền trên một chiếc máy giá trung
vị. Đổi lại, $\kappa$ giảm 100 lần và toàn bộ các thuật toán bậc một trở nên so
sánh được với nhau trong thời gian chạy hợp lý.

Đánh đổi ấy chỉ có lợi trong hai điều kiện. Thứ nhất, đường cong
cross-validation phải phẳng, vì với một đường cong có cực tiểu rõ thì khoản mất
kia lớn hơn 5 đơn vị tiền rất nhiều. Thứ hai, bài toán phải giải bằng phương
pháp lặp, vì với thuật toán bậc hai thì $\kappa$ không chi phối số vòng lặp nên
khoản được cũng biến mất.

Tổng lại, ba quyết định của chương này đưa $\kappa$ từ mức mà gradient descent
sau \SI{13.9}{\second} vẫn dừng ở sai số $10^{-2}$ xuống mức nó chạm giới hạn số
học trong \SI{8.4}{\second}, và cái giá đo được trên bài toán ứng dụng là 5 đơn
vị tiền mỗi máy. Bài toán tối ưu hóa mà các chương sau khảo sát là bài toán sau ba quyết
định ấy, với $L$, $\mu$ và $\kappa$ đã ghi ở bảng~\ref{tab:constants}.
